\section{Discussion}

\subsection{The Goldilocks Zone}

Our key finding---that moderate accommodation predicts highest engagement---challenges linear extrapolations from CAT. We propose the \textbf{accommodation-artificiality tradeoff} as one potential explanation: while some accommodation signals attentiveness, excessive accommodation may potentially trigger uncanny valley responses. However, this mechanism remains hypothetical; the inverted-U pattern could also reflect content confounds (e.g., formal topics having naturally different continuation patterns) or other unobserved factors.

\subsection{Asymmetric Bidirectional Effects}

AI-to-human accommodation ($r=0.152$) vastly exceeds human-to-AI ($r=0.013$)---an 11$\times$ difference. This suggests designers cannot rely on natural user accommodation; the burden falls on AI systems.

\subsection{Limitations}

\begin{itemize}
    \item \textbf{Correlational, not causal}: Association observed, causation unestablished
    \item \textbf{Single dataset}: Anthropic hh-rlhf only; generalization untested
    \item \textbf{Engagement proxy}: Continuation $\neq$ satisfaction
    \item \textbf{Content confounds}: Formal topics may have systematically different continuation patterns independent of accommodation
    \item \textbf{Selection effects}: hh-rlhf contains chosen vs rejected response pairs, potentially introducing selection bias
\end{itemize}
